\appendix{Представление графического материала}

Графический материал, выполненный на отдельных листах,
изображен на рисунках А.1--А.\arabic{числоПлакатов}.
\setcounter{числоПлакатов}{0}

\renewcommand{\thefigure}{А.\arabic{figure}} % шаблон номера для плакатов

\begin{landscape}

\begin{плакат}
    
\includegraphics[width=0.82\linewidth]{Сведения.eps}
    \заголовок{Сведения о ВКРБ}
    \label{pl1:image}
\end{плакат}

\begin{плакат}
    \includegraphics[width=0.82\linewidth]{Цель и задачи разработки.eps}
    \заголовок{Цель и задачи разработки}
    \label{pl2:image}
\end{плакат}

\begin{плакат}
	\includegraphics[width=0.82\linewidth]{Язык Common Lisp.eps}
	\заголовок{Язык Common Lisp}
	\label{pl3:image}
\end{плакат}

\begin{плакат}
\includegraphics[width=0.82\linewidth]{Архитектура компилятора.eps}
\заголовок{Архитектура компилятора}
\label{pl4:image}
\end{плакат}

\begin{плакат}
\includegraphics[width=0.82\linewidth]{Компиляция.eps}
\заголовок{Фаза компиляции}
\label{pl5:image}
\end{плакат}

\begin{плакат}
\includegraphics[width=0.82\linewidth]{Оптимизация.eps}
\заголовок{Фаза оптимизации}
\label{pl6:image}
\end{плакат}

\begin{плакат}
\includegraphics[width=0.82\linewidth]{Генерация.eps}
\заголовок{Фаза генерации кода}
\label{pl7:image}
\end{плакат}

\begin{плакат}
\includegraphics[width=0.82\linewidth]{Ассемблер.eps}
\заголовок{Фаза ассемблирования}
\label{pl8:image}
\end{плакат}

\begin{плакат}
\includegraphics[width=0.82\linewidth]{ВМ.eps}
\заголовок{Архитектура виртуальной машины}
\label{pl9:image}
\end{плакат}

\begin{плакат}
\includegraphics[width=0.82\linewidth]{Тесты.eps}
\заголовок{Тестирование}
\label{pl10:image}
\end{плакат}

\begin{плакат}
\includegraphics[width=0.82\linewidth]{Заключение.eps}
\заголовок{Заключение}
\label{pl11:image}
\end{плакат}

\end{landscape}
